\section{二元对称信道：从翻转概率到容量}
\label{sec:07-Information-Theory/05-Channel-Coding/02}

每个比特以概率 \(p\) 独立翻转。最省事的对策是把同一位发三遍，收端按多数表决。译码出错要求至少翻两次，

\[
P_e=3p^2(1-p)+p^3.
\]

\(p=0.1\) 时错误率从 \(0.1\) 掉到 \(0.028\)，代价是速率从每次使用 1 比特（bit）掉到 \(1/3\)。想再稳一点就发五遍、七遍，错误率继续下降，速率则一路奔向零。

按这个思路想下去，可靠性像是要拿速率来买，买多少全看你多怕出错。本节要说明这个印象不对：在同一条信道上，错误率趋于零并不需要速率趋于零，两者之间存在一条固定的分界。二元对称信道足够简单，能把这条分界精确算出来。

\subsection{模型与一张转移图}

二元对称信道 \(\mathrm{BSC}(p)\) 可以写成加噪的形式

\[
Y=X\oplus Z,\qquad Z\sim\Bernoulli(p),\ Z\perp X,
\]

其中 \(\oplus\) 是模 2 加法。等价的转移矩阵是

\[
W=\begin{pmatrix}1-p&p\\p&1-p\end{pmatrix}.
\]

\begin{center}
\begin{tikzpicture}[x=1cm,y=1cm,font=\small,>=stealth]
  \node[font=\footnotesize,black!70] at (0,2.5) {输入 \(X\)};
  \node[font=\footnotesize,black!70] at (4,2.5) {输出 \(Y\)};
  \draw[black!55] (0,1.7) circle (0.3) node {\(0\)};
  \draw[black!55] (0,0) circle (0.3) node {\(1\)};
  \draw[black!55] (4,1.7) circle (0.3) node {\(0\)};
  \draw[black!55] (4,0) circle (0.3) node {\(1\)};
  \draw[->,blue!65,thick] (0.32,1.7) -- (3.68,1.7);
  \draw[->,blue!65,thick] (0.32,0) -- (3.68,0);
  \draw[->,red!65] (0.26,1.55) -- (3.74,0.15);
  \draw[->,red!65] (0.26,0.15) -- (3.74,1.55);
  \node[above,font=\footnotesize,blue!65] at (2,1.74) {\(1-p\)};
  \node[below,font=\footnotesize,blue!65] at (2,-0.04) {\(1-p\)};
  \node[font=\footnotesize,red!65] at (1.15,0.62) {\(p\)};
  \node[font=\footnotesize,red!65] at (2.85,0.62) {\(p\)};
\end{tikzpicture}

\emph{``对称''指两条红边的概率相同：0 翻成 1 与 1 翻成 0 一样容易。这份对称后面会直接决定最优输入是公平的。}
\end{center}

给定 \(X\) 之后，\(Y\) 的全部悬念都来自噪声位 \(Z\)，所以 \(H(Y\given X)=h_2(p)\)，其中 \(h_2\) 是\hyperref[sec:07-Information-Theory/01-Information-and-Entropy/02]{``Shannon 熵与二元熵''}里那条钟形曲线。于是

\[
I(X;Y)=H(Y)-h_2(p).
\]

\subsection{最优输入不是偏好，是对称性逼出来的}

令 \(q=P(X=1)\)。输出为 1 的概率是 \(r=q(1-p)+(1-q)p\)，因此

\[
I(X;Y)=h_2\bigl(p+q(1-2p)\bigr)-h_2(p).
\]

减号右边这一项由噪声定死，输入碰不到它。要把互信息推高，只能把输出熵推高；二元变量的熵至多 1 比特，而 \(q=1/2\) 恰好让 \(r=1/2\)，上界取到。所以

\[
C_{\mathrm{BSC}(p)}=1-h_2(p)
\]

比特每次使用，由公平输入达到。这里的``公平输入最好''不是审美判断，而是转移矩阵对称的直接后果：两个输入在信道眼里地位相同，偏袒任何一方都会白白压低输出端的可用变化。信道一旦不对称，最优输入通常也就不公平了。

\subsection{最坏的翻转概率是 0.5，不是 1}

把 \(C(p)=1-h_2(p)\) 画出来，这条曲线本身就是一条判断。

\begin{center}
\begin{tikzpicture}[x=6.4cm,y=2.4cm,font=\small,>=stealth]
  \draw[->,black!55] (-0.04,0) -- (1.13,0) node[right,font=\footnotesize] {\(p\)};
  \draw[->,black!55] (0,-0.05) -- (0,1.22) node[above,font=\footnotesize] {\(C\)  (bit/use)};
  \draw[black!25,dashed] (0.5,0) -- (0.5,1.05);
  \draw[black!25,dashed] (0,1) -- (1,1);
  \draw[black!25,dashed] (0.1,0) -- (0.1,0.531) -- (0,0.531);
  \draw[blue!75,thick] plot[domain=0.0015:0.9985,samples=200]
    (\x,{1+(\x*ln(\x)+(1-\x)*ln(1-\x))/ln(2)});
  \fill[red!75] (0.5,0) circle (1.4pt);
  \fill[black!70] (0.1,0.531) circle (1.4pt);
  \fill[black!70] (0,1) circle (1.4pt);
  \fill[black!70] (1,1) circle (1.4pt);
  \node[below,font=\footnotesize] at (0,-0.06) {\(0\)};
  \node[below,font=\footnotesize] at (0.5,-0.06) {\(0.5\)};
  \node[below,font=\footnotesize] at (1,-0.06) {\(1\)};
  \node[left,font=\footnotesize] at (-0.01,1) {\(1\)};
  \node[left,font=\footnotesize] at (-0.01,0.531) {\(0.53\)};
  \node[right,font=\footnotesize,black!75] at (0.115,0.60) {\(p=0.1\)};
  \node[font=\footnotesize,red!75,align=center] at (0.5,0.22) {全噪声\\\(C=0\)};
  \node[font=\footnotesize,black!70,align=center] at (0.13,1.14) {无噪};
  \node[font=\footnotesize,black!70,align=center] at (0.87,1.14) {全翻转，仍然无损};
\end{tikzpicture}

\emph{曲线在 \(p=1/2\) 处触零，在两端各回到 1，关于 \(p=1/2\) 对称。请注意最坏的点在中间而不是右端：翻转得最多的信道并不是最差的信道。}
\end{center}

\(p=0\) 时容量为 1 不奇怪。有意思的是 \(p=1\)：每一位都确定翻转，收端照着规则再翻一次就原样恢复，容量仍是 1。一般地，\(\mathrm{BSC}(p)\) 与 \(\mathrm{BSC}(1-p)\) 只差一个确定性的输出取反，这是可逆变换，不损失任何信息，因此 \(C(p)=C(1-p)\)，通常只研究 \(0\le p\le1/2\)。

真正致命的是 \(p=1/2\)。此时输出与输入独立，收端看到的 0 和 1 与你发了什么毫无关系，容量归零。所以``错得多''并不等于``信息少''，破坏信息的是不可预测——是收端无法把输出反推回输入的那种混淆。可逆的系统性偏差可以校正，随机的对半混淆不能。这条区分在数据工作里同样管用：一个稳定的标注偏移可以事后纠回来，一批随机乱标的标签则是净损失。

\subsection{\texorpdfstring{\(0.531\) 这个数在说什么，不在说什么}{0.531 这个数在说什么}}

\(p=0.1\) 时 \(h_2(0.1)\approx0.469\)，容量约 \(0.531\) 比特每次使用。

它不表示单次发送有 \(53.1\%\) 的概率正确，也不表示存在某个短码，速率恰好 \(0.531\) 且零错误。它表示的是渐近的事：对任意 \(R<0.531\)，存在一列越来越长的编码方案，使整块消息译错的概率趋于零；而长期看，可靠速率越不过这条线。

顺带给三重重复码判个分。它的速率 \(1/3\approx0.333\)，离 \(0.531\) 还差一大截，可它的错误率也远没到``任意小''。差距的来源不是冗余不够，是冗余的花法太笨：逐位重复只在一维上做文章，而好的码把大量消息摊到高维二元空间里，让码字彼此隔开足够的汉明距离，于是典型的噪声图样仍不足以把一个码字推到另一个码字的地盘。这套``用距离换纠错能力''的算账方式在\hyperref[sec:04-Discrete-Mathematics/05-Coding-and-Polynomial-Methods/02]{``码距决定检测与纠错能力''}里有代数版本，两边说的是同一件事的不同侧面。

\subsection{什么时候不能直接套这条公式}

\(1-h_2(p)\) 的适用面比它的名气窄。它要求错误独立、同概率、并且已经被压成硬判决的 0/1。现实中常见的偏离有几种：0 到 1 与 1 到 0 的概率不等，此时信道不再对称，最优输入也不再公平；错误成串出现而非独立，信道有了记忆；信道状态随时间变化，或只有收端知道当前状态；编码器还受输入代价或组成约束的限制。

最容易被忽略的是最后一种情形：接收器本来握有连续幅度，却被提前量化成 0/1。这一步硬判决本身是对 \(Y\) 的后处理，按数据处理不等式只会丢信息，于是算出来的 \(1-h_2(p)\) 只是原信道可达能力的一个下界。实际系统之所以坚持把软信息（对数似然比）一路送进译码器，正是为了不提前扔掉这部分。

下一节把视角从一条具体信道抬到一般情形：容量为什么可以由一个单字母的优化问题定义，而这个定义又凭什么管得住长块编码的命运。
